\scanpage{12}\mspage{11}

L'isomorphisme de la fin de d) nous fournit une section canonique de

\[
\underline{\mathrm{Pic}}_{B/S}^{-\delta'}\overset A\times
(\underline{\mathrm{Pic}}^\delta_{A/S}\overset B\times(\deleted{\unsure{B}}A,-\widetilde{\delta'}))
\]

où $\delta'=-\varphi(\delta)$ comme dessus. En d'autres termes, on trouve un isomorphisme canonique

\[
\boxed{\underline{\mathrm{Pic}}^{\delta'}_{B/S}\simeq\underline{\mathrm{Pic}}^{\deleted{\unsure{=|-}}-\delta}_{A/S}\overset B\times(A,\widetilde{\delta'})}\tag*{(*)}
\]

\margin{$ {}_2 A$}

% slanted
\margin{\uline{NB} Comme $\underline{\mathrm{Pic}}^{-\delta}_{A/S}\simeq\underline{\mathrm{Pic}}^\delta_{A/S}$ (isom canonique) on peut supprimer si on veut les signes dans notre formule\ldots{} Prouver qu'après passage aux carrés des torseurs, on trouve un isom qui est compatible avec les sections canoniques, donc que $(*)$ provient d'un isom entre les torseurs \unsure{correspondants} au Groupe $ {}_2 A$\ldots. \uline{à vérifier}}

\uline{Exemple} Supposons que $\delta$ soit de degré 1, \deleted{\unsure{$\delta'$ est u}} $\widetilde\delta:A\xrightarrow{\sim}B$, et $\widetilde{\delta'}:B\simeq A$ est l'inverse de $\widetilde\delta$. Alors $\delta'$ \deleted{et \unsure{$\delta'$ c'est la}} \inserted{est la} polarisation déduite de $\delta$ par transport de structure à l'aide de $\widetilde\delta$. On trouve un isom., \unsure{grâce|via} à $\widetilde\delta$, $\underline{\mathrm{Pic}}^{\delta'}_{B/S}\simeq\underline{\mathrm{Pic}}^\delta_{A/S}$ (compatible \deleted{\unsure{avec $\underline{\mathrm{Pic}}$}} avec $\deleted{\unsure{B}}A\xrightarrow[\sim]{\widetilde\delta}B$) et l'isom considéré peut s'interpréter comme un isom \inserted{de $B$-torseurs}

\[
\underline{\mathrm{Pic}}^\delta_{A/S}\simeq\underline{\mathrm{Pic}}^{-\delta}_{A/S}
\]

\deleted{C'est là l'isomorphisme canonique bien connu, défini par la condition de \unsure{transformer} le diviseur $\Theta^{\unsure{-|{}}}$ en $\Theta^{-}$.} \deleted{\unsure{sur $A\times B$|$A\times B$}} \deleted{\unsure{Je crois que c'est la propriété}\illegible{2}\ldots}

\uline{Remarque} La section du torseur de la fin de d) est stable par autom. de $(X/S,\delta)$ ; dans le cas actuel, on peut prendre les autom par translations dans $X=A$. On trouve que la section par $a$ opère sur

\[
\underline{\mathrm{Alb}}_{X/S}^{\chi(\delta)}\times\underline{\mathrm{Pic}}^{\varphi(\delta)}_{B/S}
\overset A\times (P^\delta\overset B\times(A,\widetilde{\varphi(\delta)}))
\simeq\underline{\mathrm{Pic}}^{\varphi(\delta)}_{B/S}\overset A\times
(P^\delta\overset B\times(A,\widetilde{\varphi(\delta)}))
\]

% slanted
\margin{Attention au signe !?!}

par translation par $\chi(\delta)a+\widetilde{\varphi(\delta)}\,\widetilde\delta(a)\deleted{\unsure{a}}$ ; donc on a $[\chi(\delta)\mathrm{id}_A+\widetilde{\varphi(\delta)}\,\widetilde\delta]a=0$ pour tt $a$, d'où \deleted{\unsure{résulte}} $\widetilde{\varphi(\delta)}\,\widetilde\delta=-\chi(\delta)\mathrm{id}_A$
